\documentclass[twocolumn]{iopjournal}

\usepackage{amsmath,amssymb,amsfonts}

\fancyhfoffset[L]{0mm}
\setlength{\marginparwidth}{0pt}
\setlength{\marginparsep}{0pt}

% ==========================================
% ZENODO PREPRINT HEADERS (ACTIVE)
% ==========================================
\renewcommand{\articletype}[1]{%
  {\vspace*{-8mm}\noindent \Large \sf Preprint}\par
  \vspace*{6mm}\noindent{\scriptsize \sf{\bfseries \MakeUppercase{#1}}}}

\fancyhead[L]{}
\fancyhead[R]{}

\begin{document}

\articletype{Supplementary Material}

\title{Supplementary Material: The Constitutive Law\\
from Exchange Symmetry}

\author{Stephen Nelson}

\affil{Field Dynamics Research}

\ead{stephen@fielddynamics.org}

\begin{abstract}
The tetrad equation produces two parity-related bounded regions,
TetA and TetB, proved in the main
letter~\cite{Nelson2026geometry}. Any constitutive response
$\mu(x)$ measuring the TetA fraction at acceleration $x$ must
satisfy a duality: exchanging TetA and TetB sends $x \to 1/x$
and $\mu \to 1 - \mu$. Defining the odds ratio
$\Omega = \mu/(1 - \mu)$, this duality becomes
$\Omega(1/x) = 1/\Omega(x)$. The map $\Omega$ must be rational
(to preserve the tetrad equation's algebraic structure) and
bijective (for deterministic decoding); a rational bijection
is necessarily fractional linear~\cite{Hartshorne1977}. Within
this class, the exchange symmetry and the endpoint condition
$\mu(0) = 0$ together force $\Omega(x) = x$, from which
$\mu(x) = x/(x+1)$ follows by algebraic normalisation. The
midpoint value $\mu(1) = 1/2$ is not assumed; it is a derived
consequence. This derivation uses no coupling polynomial, no
rank counting, no field equation, and no equilibrium postulates.
It requires only the exchange symmetry between the two
parity-related roots of the tetrad equation.
\end{abstract}

% ----------------------------------------------------------------------

\section*{Exchange Symmetry}

The main letter~\cite{Nelson2026geometry} and its
supplement~\cite{Nelson2026scheme} prove that the tetrad
equation $g_{\mu\nu} = \eta_{ab}\, e^a{}_\mu\, e^b{}_\nu$
produces exactly two bounded tetrahedral regions in its spatial
sector, related by spatial parity $P: e^i \to -e^i$. The
identity component gives TetA; the parity component gives TetB.
Their union is the stellated octahedron.

Let $x = g/a_0$ be the dimensionless acceleration ratio. Define
$\mu(x) \in [0,1]$ as the fraction of the total field capacity
carried by TetA at acceleration $x$, so that $1 - \mu(x)$ is
the fraction carried by TetB.

Spatial parity exchanges TetA and TetB. Under this exchange,
the acceleration ratio inverts ($x \to 1/x$) and the two
fractions swap. The constitutive response must therefore satisfy
the duality condition:
\begin{equation}
  \mu(1/x) = 1 - \mu(x).
  \label{eq:duality}
\end{equation}
This is the exchange symmetry of the stellated octahedron
expressed as a functional equation on $\mu$.

% ----------------------------------------------------------------------

\section*{Endpoint Conditions}

Two endpoint conditions follow from the definitions of the
TetA and TetB fractions:

\paragraph{$\mu(0) = 0$.} At zero acceleration, the TetA
fraction vanishes. The field sits entirely in TetB.

\paragraph{$\mu(x) \to 1$ as $x \to \infty$.} At arbitrarily
large acceleration, the TetA fraction saturates.

These are structural definitions: $\mu = 0$ means TetB alone,
$\mu = 1$ means TetA alone.

% ----------------------------------------------------------------------

\section*{The Odds Ratio}

Define the odds ratio between TetA and TetB:
\begin{equation}
  \Omega(x) = \frac{\mu(x)}{1 - \mu(x)}.
  \label{eq:odds}
\end{equation}
This measures how many times more capacity TetA carries than
TetB at acceleration $x$. Under the exchange symmetry
Eq.~\eqref{eq:duality}:
\begin{align}
  \Omega(1/x)
  &= \frac{\mu(1/x)}{1 - \mu(1/x)}
   = \frac{1 - \mu(x)}{\mu(x)} \notag \\
  &= \frac{1}{\Omega(x)}.
  \label{eq:odds_duality}
\end{align}
The odds ratio inverts when the acceleration inverts. This is the
exchange symmetry in its simplest algebraic form.

The endpoint conditions translate to:
\begin{equation}
  \Omega(0) = 0, \qquad \Omega(x) \to \infty
  \text{ as } x \to \infty.
  \label{eq:odds_endpoints}
\end{equation}

% ----------------------------------------------------------------------

\section*{Rational Bijectivity}

The map $\Omega(x)$ must be rational, to preserve the algebraic
structure of the tetrad equation under composition, and bijective,
so that a single acceleration $x$ corresponds to exactly one odds
ratio and vice versa. A rational bijection on the projective line
is necessarily fractional linear; this is a classification theorem
of projective geometry~\cite{Hartshorne1977}. The general form is:
\begin{equation}
  \Omega(x) = \frac{ax + b}{cx + d}.
  \label{eq:mobius_general}
\end{equation}

% ----------------------------------------------------------------------

\section*{The Unique Map}

\paragraph{Proposition.} \textit{Let $\Omega(x) = (ax+b)/(cx+d)$
be a rational bijection satisfying $\Omega(1/x) = 1/\Omega(x)$
and $\Omega(0) = 0$. Then $\Omega(x) = x$.}

\paragraph{Proof.}
The endpoint condition $\Omega(0) = b/d = 0$ requires $b = 0$.
The map reduces to $\Omega(x) = ax/(cx + d)$.

Substituting into the exchange condition
$\Omega(1/x) = 1/\Omega(x)$:
\begin{equation}
  \frac{a}{c + dx} = \frac{cx + d}{ax}.
  \label{eq:exchange_sub}
\end{equation}
Cross-multiplying: $a^2 x = (c + dx)(cx + d)$. Expanding
the right side and comparing coefficients of $x^2$, $x^1$,
and $x^0$:
\begin{align}
  x^2: \quad & 0 = cd, \label{eq:x2} \\
  x^1: \quad & a^2 = c^2 + d^2, \label{eq:x1} \\
  x^0: \quad & 0 = cd. \label{eq:x0}
\end{align}
From Eq.~\eqref{eq:x2}, either $c = 0$ or $d = 0$.

If $d = 0$: $\Omega(x) = ax/(cx) = a/c$, a constant. This is
not bijective. Rejected.

If $c = 0$: $\Omega(x) = ax/d$, and Eq.~\eqref{eq:x1} gives
$a^2 = d^2$, so $a = \pm d$. Since $\Omega(x) > 0$ for $x > 0$,
we require $a/d > 0$, giving $a = d$. Therefore:
\begin{equation}
  \boxed{\Omega(x) = x.}
  \label{eq:odds_result}
\end{equation}

The odds ratio between TetA and TetB is the acceleration ratio
itself. This is not assumed. It is the unique rational bijection
satisfying the exchange symmetry and the endpoint condition.
\hfill $\square$

% ----------------------------------------------------------------------

\section*{The Constitutive Law}

From $\Omega(x) = x$, inverting the definition
Eq.~\eqref{eq:odds}:
\begin{equation}
  \mu = \frac{\Omega}{1 + \Omega} = \frac{x}{1 + x}.
  \label{eq:mu}
\end{equation}
\begin{equation}
  \boxed{\mu(x) = \frac{x}{x + 1}.}
\end{equation}

\paragraph{Derived consequences (not inputs).}
The midpoint value $\mu(1) = 1/2$ was not assumed. It follows:
\begin{equation*}
  \mu(1) = \frac{1}{1 + 1} = \frac{1}{2}.
\end{equation*}
At the parity midpoint, TetA and TetB contribute equally. This
is a consequence of the exchange symmetry, not an independent
condition.

The complement symmetry $\mu(x) + \mu(1/x) = 1$ was not
assumed. It is the duality condition Eq.~\eqref{eq:duality},
which was the starting point.

The screening identity $x(1 - \mu) = \mu$ is not imposed. It
is the odds-ratio identification $\mu/(1-\mu) = x$ rearranged:
multiply both sides by $(1-\mu)$ to obtain $\mu = x(1-\mu)$.
This is an algebraic identity of Eq.~\eqref{eq:mu}.

% ----------------------------------------------------------------------

\section*{Relation to the Other Routes}

The main paper derives $\mu(x) = x/(x+1)$ by two additional
routes:

\vspace{6pt}
\noindent
\begin{tabular*}{\columnwidth}{@{\extracolsep{\fill}}lp{0.37\columnwidth}@{}}
\hline
Route & Method \\
\hline
1 (companion note) & Graded rank counting of the tetrad
  equation yields coupling polynomial; parity interpolation
  and rational normalisation produce
  $\mu$~\cite{Nelson2026metric} \\
2 (companion note) & AQUAL field equation, three physical
  postulates, single-channel detailed
  balance~\cite{Nelson2026equilibrium} \\
3 (this note) & Exchange symmetry of the two parity-related
  roots; odds ratio identified with acceleration ratio by
  rational bijectivity \\
\hline
\end{tabular*}
\vspace{6pt}

Route~1 reads the algebraic content of the tetrad equation.
Route~2 reads the field's equilibrium response. Route~3 (this
note) reads the geometry's duality. No two routes share their
method:

\begin{itemize}
\item Route~1 uses the coupling polynomial. This note does not.
\item Route~2 uses the AQUAL field equation. This note does not.
\item This note uses the odds ratio. Routes~1 and~2 do not.
\end{itemize}

All three arrive at the same function. The main letter presents
the convergence of the three routes.

% ----------------------------------------------------------------------

\section*{Summary}

The tetrad equation has two parity-related spatial roots. Their
exchange symmetry requires $\mu(1/x) = 1 - \mu(x)$. Defining
the odds ratio $\Omega = \mu/(1-\mu)$, this becomes
$\Omega(1/x) = 1/\Omega(x)$. Rational bijectivity restricts
$\Omega$ to the fractional linear class. The endpoint condition
$\Omega(0) = 0$ sets the numerator constant to zero. The
exchange identity then forces $cd = 0$ and $a^2 = d^2$, with
the bijective branch giving $c = 0$ and $a = d$. The result is:
\begin{equation}
  \Omega(x) = x, \qquad \mu(x) = \frac{x}{x + 1}.
\end{equation}
The odds ratio between the two tetrahedral sectors is the
acceleration ratio itself. The constitutive law is its
normalisation. The midpoint value $\mu(1) = 1/2$ is derived,
not assumed.

This function was identified empirically by
Milgrom~\cite{Milgrom1983} and selected by fitting by
Famaey and Binney~\cite{FamaeyBinney2005}. The present
derivation obtains it from the exchange symmetry of the
two parity-related roots of the equation that defines the
gravitational field.

% ----------------------------------------------------------------------

\begin{thebibliography}{99}

\bibitem{Nelson2026geometry}
S.~Nelson,
\textit{The Coordinate-Space Geometry of the Tetrad Frame},
Class.\ Quantum Grav.\ (2026), submitted.

\bibitem{Nelson2026scheme}
S.~Nelson,
\textit{Supplementary Material: The Tetrad Scheme and the Two
Bounded Regions} (2026).

\bibitem{Nelson2026metric}
S.~Nelson,
\textit{Supplementary Material: The Constitutive Law Read from
the Metric} (2026).

\bibitem{Nelson2026equilibrium}
S.~Nelson,
\textit{Supplementary Material: The Constitutive Law from
Equilibrium Field Closure} (2026).

\bibitem{Hartshorne1977}
R.~Hartshorne,
\textit{Algebraic Geometry}
(Springer, New York, 1977).

\bibitem{Milgrom1983}
M.~Milgrom,
\textit{A modification of the Newtonian dynamics as a possible
alternative to the hidden mass hypothesis},
Astrophys.\ J.\ \textbf{270} (1983) 365.

\bibitem{FamaeyBinney2005}
B.~Famaey and J.~Binney,
\textit{Modified Newtonian dynamics in the Milky Way},
Mon.\ Not.\ R.\ Astron.\ Soc.\ \textbf{363} (2005) 603.

\end{thebibliography}

\end{document}
